% ************************** Thesis Summary *****************************

\begin{summary}
\addcontentsline{toc}{section}{\textbf{Summary}}

\subsection*{Motivation}
\addcontentsline{toc}{subsection}{Motivation}
It is plain to see that economic decision-making is a fundamental aspect of modern human existence. Even before sitting down at their desk, a commuter must integrate multiple decision variables to calculate which mode of transport to take, is bombarded with advertisement for consumer goods, and faces a choice of breakfast options that would be unimaginable even a millennium ago. The brain has the unenviable duty to gather required information, incorporate it into a common framework, and decide between alternatives for each of these tasks.

Over the last 500 years, economics has found success in explaining what people decide when they are deciding. Utility has emerged as the fundamental variable of interest which is maximised along economic choices. These subjective values present in any choice are generated internally to the individual faced with alternatives and must be elicited via careful study of behaviour. Over the past 50 years, auction theory has proved successful in developing mechanisms which incentivise participants to declare the utilities of offered goods on a single trial basis.

Theories of utility have proved successful both prescriptively, in setting axioms for how a hyper-rational decision maker should act, and descriptively when studying how and when humans break these 'rules' (\cite{luce2000utility}). As impressive as the economics literature is, it cannot tell us how the brain manages to navigate these choices every single day of our lives\footnote{Neither literally how neurons calculate and compare utilities, nor philosophically how a 1500g piece of meat manages to hold its own against intelligently designed computer programs. As an essay beloved by the author puts it (\cite{bisson1991theyre}):

``We probed [humans]. They're meat all the way through.''\\
``[They don't have a brain?]''\\
``Oh, there is a brain all right. It's just that the brain is made out of meat!''\\
``So... what does the thinking?''\\
``You're not understanding, are you? The brain does the thinking. The meat.''\\
``Thinking meat! You're asking me to believe in thinking meat!''}.


More recently, neuroscience has identified the activity of dopaminergic cells in the midbrain as a putative substrate for utility. These neurons have been shown to encode and integrate multiple variables crucial to this function such as the magnitude, probability, and temporal delay of rewards. By combining these parameters into a common signal, dopaminergic neurons facilitate the direct comparison of and selection from choices that encompass a vast decision space. By reinforcing choices that maximise utility, these neurons serve a crucial role in balancing an organism’s motivations and disturbances to this system have been implicated in severe neuropsychiatric disorders such as Parkinson’s disease and schizophrenia.

Central to research into the function of midbrain dopaminergic neurons are electrophysiological recording methods, which allow researchers to study the behaviour of these cells with unparalleled temporal and spatial resolution. Such recordings require direct access to nuclei nested deep within the brain and so are unsuitable for studies of human decision-making. Instead, most cellular studies of economic decision-making have relied on rodent or non-human primate models, limiting the sophistication of the choice tasks that can be employed.

The motivation for this thesis is to overcome these limitations so that we might gain a greater understanding of the dopaminergic utility signal. We aim to investigate the activity of midbrain dopaminergic neurons in rhesus macaques trained on a task inspired by auction theory to provide a real time readout of utility while observing changes in neuronal activity.

\subsection*{Organisation of the Thesis}
\addcontentsline{toc}{subsection}{Organisation of the Thesis}

The first chapter is spent reviewing the fields of economics and neuroscience. We chronicle the development of theories of utility, before examining the evidence that the brain, and especially the dopaminergic midbrain circuitry, encodes utility. We also discuss the experimental task paradigms that animal research into utility has employed.

In Chapter 2, we compare these paradigms to a second-price auction developed by Becker et al., 1964 and explore the theory behind what makes these ‘Becker-DeGroot-Marshak’ (BDM\nomenclature{BDM}{Becker-DeGroot-Marshak (auction)}) auctions appealing to behavioural economists and show theoretical results of a rational bidder. We conclude the chapter by looking at the growing neuroscientific literature of human economic decision-making using auction tasks, and explain why we consider the BDM an appealing choice for investigating the function of neuronal circuits of reward. 

In Chapter 3 we detail a BDM auction task we developed and trained two rhesus macaques on. The behaviour of the primates on the task is explored, and the monkey's manipulation of tasks inputs and bidding behaviour is studied. The task behaviour is also compared to the same task limited to just one auctioned good and to more traditional binary choice tasks. We show convincing evidence that the animals are able to perform the BDM auction task \textit{as if} they are maximising some internal utility.

Chapter 4 concerns the electrophysiological recording of the dopaminergic cells of the midbrain of the two trained rhesus macaques. It begins with the single-cell recording procedure and the confirmation of the dopaminergic nature of the recorded cells. We then analyse the activity of dopaminergic cells as the animal responds to the auction task and makes their bid. We show that on a trial by trial level, dopaminergic midbrain cells in the animals encode various trial variables such as auction offer, the opposing computer bid, and whether or not the monkey won the trial. Crucially, we show that the activity of dopaminergic cells in the epoch where a good is offered is correlated with the subsequent monkey's bid.

Finally, we conclude by reviewing the results presented and the potential future directions that this work implies. We also briefly discuss the current state of primate/dopaminergic research and how neuroscientific studies can move beyond staid hypotheses.
\end{summary}
